\documentclass[aps,prx,article,twocolumn,amsmath,amssymb,superscriptaddress,longbibliography]{revtex4-2}
\usepackage{graphicx}  
\usepackage{dcolumn}   
\usepackage{bm}        
\usepackage{amssymb}   
\usepackage{amsmath}
\usepackage{mathrsfs}
\usepackage{epigraph}
\usepackage{braket}
\usepackage{gensymb}
\usepackage{lmodern}
\usepackage{ tipa }
\usepackage{bbold}
\usepackage{esint}
\usepackage{mathdots}
\usepackage{xcolor}
\usepackage{appendix}
\usepackage{multirow}
\usepackage{float}
\usepackage{physics}
\usepackage[normalem]{ulem}
\usepackage{times}

\usepackage{comment}
\usepackage[colorlinks=true, linkcolor=blue, urlcolor=blue, citecolor=blue]{hyperref}


% \usepackage{lineno}
% \linenumbers


\begin{document}
\title{Anharmonic Superlattice Distortions and Charge Order\\ in Rhombohedral Multi-Layer Graphene} 	
\author{Haoran Yan}
\affiliation{Department of Chemistry, Emory University, Atlanta, GA 30322, USA}
\author{Wangqian Miao}
\affiliation{Department of Physics, The Pennsylvania State University, University Park, PA 16802, USA}
\affiliation{Department of Chemistry, Emory University, Atlanta, GA 30322, USA}
\author{Ziyan Zhu}
\affiliation{Department of Physics, Boston College, Chestnut Hill, MA, USA}
\author{Jiang-Xiazi Lin}
\email[Correspondence should be addressed to \href{mailto:jiangxiazi.lin@emory.edu}{jiangxiazi.lin@emory.edu} and \href{mailto:yao.wang@emory.edu}{yao.wang@emory.edu}]{}
\affiliation{Department of Physics, Emory University, Atlanta, GA 30322, USA}
\author{Yao Wang}
\email[Correspondence should be addressed to \href{mailto:jiangxiazi.lin@emory.edu}{jiangxiazi.lin@emory.edu} and \href{mailto:yao.wang@emory.edu}{yao.wang@emory.edu}]{}
\affiliation{Department of Chemistry, Emory University, Atlanta, GA 30322, USA}
\date{\today}
\begin{abstract} 
The flat-band electronic system without moiré potential in rhombohedral multilayer graphenes (RMG) has attracted tremendous attention because of emergent phases such as fractional Chern insulating insulators and chiral superconductors. These exotic states are widely believed to originate from strong correlations in the flat-band systems. At the same time, additional phases associated with spatial-symmetry breaking have been observed in certain regions of the phase diagram. Here, we propose a theoretical framework to explain these spatially anisotropic phases together with the observed doping hysteresis. Using the rhombohedral hexalayer graphene as an example, we show that Fermi-surface nesting within a nonmagnetic regime of the parameter space can drive a charge instability. Once established, the resulting charge-density-wave order breaks both translational and rotational symmetries and couples to anharmonic superlattice distortions, thereby accounting for the recently observed hysteretic and antimagnetic behavior. Our framework further suggests the presence of a competing subleading charge instability and strong electron-lattice coupling in RMG, which may contribute to the other correlated phases observed in the family of materials.
\end{abstract}

\maketitle

\section{Introduction}

Flat-band two-dimensional systems have become a central platform for exploring tunable strongly correlated states, because the suppression of kinetic energy amplifies the effects of electron-electron interactions\,\cite{bistritzer2011moire, nuckolls2024microscopic, andrei2021marvels}. Beyond twisted moiré systems, in which many exotic phases such as FCI and superconductivity have been observed\,\cite{cao2018unconventional,cao2018correlated,balents2020superconductivity,ju2024fractional,bernevig2025fractional}, rhombohedral multilayer graphene (RMG) has recently emerged as an especially clean and versatile alternative, in where crystalline stacking naturally produces flatbands at low doping, while displacement fields can continuously tune the band structure and geometry. Experiments have now revealed a broad range of correlated and topological phases in RMG, including spin- and valley-polarized quarter-metal states\,\cite{zhou2021half, han2023orbital, auerbach2025isospin, shi2020electronic, xie2026magnetic, aronson2025displacement}, fractional Chern insulating phases\,\cite{lu2024fractional, lu2025extended, xie2025tunable, zhang2025moire, xiao2025thickness}, and unconventional superconductivity with possible chiral pairing\,\cite{zhou2021superconductivity, han2025signatures, han2026evidence, choi2025superconductivity, holleis2025fluctuating}.

In addition to these exotic magnetic and superconducting states, recent experiments have revealed symmetry-broken stripe-like electronic orders in rhombohedral graphene, manifested as spatially anisotropic resistivity in angle-resolved transport measurements\,\cite{qin2025extreme,morissette2025superconductivity}. Intriguingly, these stripe orders are found to coexist with superconductivity and may even enhance it, raising the possibility that the anisotropy and the pairing mechanism are linked by a common underlying origin. These discussions can be further traced back to earlier studies of rhombohedral trilayer graphene, where SNOM experiments have revealed electron-phonon coupling\,\cite{zan2024electron} and mean-field simulations have suggested phonon-mediated topological superconductivity\,\cite{chou2021acoustic}. Together, these observations indicate a potential connection among anisotropic stripes, phonons, and superconductivity.

A more recent experiment has provided a more systematic investigation of the nature of this stripe order and its interplay with magnetic orders\,\cite{zhao2026hysteretical}. The stripe order is found to substantially compete with spin- and valley-polarized states, which are triggered in regimes of the phase diagram by orbital ferromagnetism, and be suppressed by magnetic field. More importantly, the anisotropic state is consistently accompanied by a doping-driven hysteresis. This hysteretic behavior closely parallels a recent observation in the quasi-one-dimensional quantum spin Hall insulator TaIrTe$_4$, where it has been attributed to strong electron--phonon coupling based on scanning-probe and Raman evidence\,\cite{tang2026bistable}. These recent observations suggest a unified picture in which anisotropic stripe orders are induced by strong coupling to anharmonic phonons.

Here we develop a minimal phenomenological framework for understanding the anisotropy and hysteresis recently observed in rhombohedral graphene. Our central idea is that, near a doping-tuned van Hove singularity, the Fermi surface becomes nested at specific wavevectors, enhancing the charge susceptibility and driving a stripe-like CDW instability that breaks both translational and rotational symmetry. In flat-band systems, where electron--phonon coupling is enhanced, the onset of this stripe order can further induce a substantial lattice distortion toward a lower-symmetry structure with a superlattice commensurate with the CDW order. Because the lattice structure in RMG is itself unstable, this distortion can be trapped in a local metastable configuration by lattice anharmonicity, allowing the structural distortion to survive even when the electronic order is no longer present at a different doping level. This mechanism naturally gives rise to resistivity hysteresis along the doping axis, and it is consistent with the strong sensitivity of the effect to magnetic field and spin polarization. Although the present framework does not explicitly include strong electronic correlations, it nevertheless provides a unified explanation for the recently observed anisotropic, nonmagnetic, and hysteretic states.


\section{Tight-Binding Model for Rhombohedral Multi-Layer Graphene}

To model the electronic structure of RMG, we employ a tight-binding description based on the underlying atomic geometry, an approach that has been widely used in parameterizing 2D materials\,\cite{fang2015ab,koshino2018maximally,miao2023truncated, vafek2023continuum}. Within this framework, the electronic hopping parameters between different C $p_z$ orbitals are approximated through the Slater--Koster (SK) parametrization, whose explicit geometry-dependent form naturally incorporates lattice distortions through the modified atomic configuration\,\cite{slater1954simplified}. The hopping amplitude between two sites separated by $\mathbf{r}$ is written as
\begin{align}
   t(\mathbf{r}) = V_{pp\pi}(r)\left(1-\frac{z^2}{r^2}\right) + V_{pp\sigma}(r)\frac{z^2}{r^2}, 
\end{align}
where $z$ is the out-of-plane component of $\mathbf{r}$ and $r = |\mathbf{r}|$. The distance-dependent $\pi$- and $\sigma$-type hopping integrals are given by
\begin{align}
V_{pp\pi}(r) &= V_{pp\pi}^0 \, e^{q_\pi (1 - r/a_\pi)} f_c(r),\\
V_{pp\sigma}(r) &= V_{pp\sigma}^0 \, e^{q_\sigma (1 - r/a_\sigma)} f_c(r),
\end{align}
with cutoff function $f_c(r) = \left[1 + e^{(r - r_c)/l_c}\right]^{-1}.$
For the calculations presented here, we use $V_{pp\pi}^0 = -2.81~\mathrm{eV}$ and $V_{pp\sigma}^0 = 0.48~\mathrm{eV}$ for the overall $\pi$- and $\sigma$-type hopping strengths, together with reference bond lengths are $a_\pi = 1.418~\text{\AA}$ and $a_\sigma = 3.349~\text{\AA}$\,\cite{herzogarbeitman2024moire}. The inclusion of long-range hopping breaks particle-hole symmetry in the model, consistent with the known electronic structure of RMG\,\cite{huang2023spin, koh2024correlated, herzogarbeitman2024moire,wolf2024magnetism,parramartinez2025nematic}.


To depict the low-energy electronic states of RMG, we incorporate a perpendicular displacement field as a layer-dependent onsite potential. This field breaks inversion symmetry and generates an energy offset between different graphene layers\,\cite{herzogarbeitman2024moire}, leading to the second-quantized term
\begin{align}
\mathcal{H}_D = \sum_{l,\alpha} \frac{|e| d_0 D}{\epsilon_r} \left(l - \frac{n-1}{2}\right) c^\dagger_{l\alpha} c_{l\alpha},
\end{align}
where $l$ labels the layer index for an $n$-layer system, $D$ denotes the displacement field, and $d_0$ is the interlayer spacing. This term induces a tunable interlayer asymmetry and opens a gap at charge neutrality.

In this paper, we focus on the rhombohedral six-layer graphene, motivated by recent angle-resolved transport experiments\,\cite{qin2025extreme}, although the model and framework apply to general systems. Figure~\ref{fig: Figure1_flatband_and_dos} presents the band structure and density of states (DOS) calculated from our model. At zero displacement field ($D=0$), the system exhibits nearly gapless bands, with a characteristic flat-band regime near the $K$ and $K'$ point (highlighted by bold lines). Physically, these low-energy electronic states in the flat-band regime are predominantly localized on the outermost layers \cite{zhang2011spontaneous,lau2021designing, zhang2024correlated, xiao2025thickness}, which makes them especially sensitive to an applied displacement field compared with higher-energy states. Once a finite displacement field is introduced, a gap opens and the top valence band develops a Mexican-hat-like dispersion together with Fermi-surface pockets (see the right panel in Fig~\ref{fig: Figure1_flatband_and_dos})\,\cite{wolf2024magnetism,parramartinez2025nematic}. 

\begin{figure} [!t]
\centering 
\includegraphics[width=8.5cm]{Figure1_flatband_and_dos_v2.eps} \vspace{-3mm}
\caption{Band structures of rhombohedral hexalayer graphene near valley $K$ for zero electric displacement field (left) and for a finite displacement field of $D=-160\,\text{mV/nm}$ (right). The tunable flat-band regime is highlighted by bold lines. The corresponding DOS for each case is shown in the right panel, where the van Hove singularity relevant to this paper is marked.}
\label{fig: Figure1_flatband_and_dos} 
\end{figure}

This flat-band sector is widely believed to underlie the observed unconventional superconductivity\,~\cite{chou2025intravalley,dong2026superconductivity,gaggioli2025spontaneous,li2025berry,geier2025chiral,parramartinez2025band,shavit2025ephemeral,jahin2026enhanced} and possible Wigner-crystal-like behavior as primarily associated with the flat-band electronic states\,~\cite{dong2024anomalous,dong2024theory,dong2024stability,tan2024parent,zeng2024sublattice,kwan2025moire,miao2026various}, whereas other phenomena, such as orbital ferromagnetism and fractional quantum Hall effects, may also involve contributions from additional bands\,\cite{huang2023spin,wolf2024magnetism,guo2024fractional,huang2025fractional,yu2025moire}. A key consequence of the flat low-energy dispersion is the emergence of a van Hove singularity (VHS) near the charge neutral point (CNP), which will be the primary focus of this paper. It is worth noting that the band tops of lower bands generate VHSs as well, although with lower DOS, and these features will play a role in the later discussion. 




\section{Two-Order Phenomenological Model for Hysteresical Anisotropy}

Near the VHS regime, the strongly enhanced DOS makes the system particularly susceptible to ordering instabilities driven by electronic interactions or electron-phonon coupling. At the same time, the nature of these instabilities is strongly influenced by the geometry of the Fermi surface. When the system is hole doped into the VHS regime, the Fermi surface reconstructs into triangular pockets, as shown in Fig.~\ref{fig: Figure2_FS_nesting}(b). In this geometry, different segments of the pockets become locally parallel and are connected by characteristic finite wavevectors $\mathbf q_{\mathrm{nesting}}$, highlighted in the zoomed-in panel. This Fermi-surface geometry naturally promotes nesting between distinct low-energy regions, thereby favoring a CDW-like ordering instability near the VHS regime.

\begin{figure}[!b]
\centering
\includegraphics[width=8.5cm]{Figure2_FS_nesting_v2.eps}\vspace{-3mm}
\caption{Fermi surfaces of rhombohedral hexalayer graphene near the $K$ valley at displacement field $D=-160$\,nV/nm for hole doping of (a) $n_H = 1.57 \times 10^{12}\,\mathrm{cm}^{-2}$, far from the VHS, and (b) $n_H = 0.117 \times 10^{12}\,\mathrm{cm}^{-2}$, near the VHS. The zoomed-in panel illustrates the approximate nesting condition of the near-VHS Fermi surface at wavevector $\mathbf{q}_{\mathrm{nesting}}$.
}
\label{fig: Figure2_FS_nesting}
\end{figure}

To quantify this instability, we evaluate the static charge susceptibility within the generalized Lindhard formalism in the non-interacting limit,
\begin{align}\label{eq:susceptibility}
\chi(\mathbf{q}) =
\sum_{\mathbf{k},m,n}
|\left\langle u_{m,\mathbf k+\mathbf q} \middle| u_{n,\mathbf k}\right\rangle|^2
\frac{f(\epsilon_{m,\mathbf{k}})-f(\epsilon_{n,\mathbf{k}+\mathbf{q}})}
{\epsilon_{n,\mathbf{k}+\mathbf{q}}-\epsilon_{m,\mathbf{k}}+i\eta},
\end{align}
where $m$ and $n$ label band indices, $f(\epsilon)$ is the Fermi--Dirac distribution function, and $\eta$ is a small \textit{ad hoc} broadening parameter introduced to regularize the susceptibility. The form factor $\left\langle u_{m,\mathbf k+\mathbf q} \middle| u_{n,\mathbf k}\right\rangle$ describes the overlap between eigenstates in the corresponding bands and captures the important quantum-geometric effects present in RMG\,\cite{parramartinez2025band, jahin2026enhanced}. 

Based on the susceptibility, we formulate an effective Ginzburg--Landau theory for the electronic order and its associated phase transition. For simplicity, we retain only the charge order associated with the dominant nesting wavevector $\mathbf q_{\mathrm{nesting}}$, represented by the order parameter $\psi_{\vb*{q}}$. Although the microscopic details of the electronic interactions and wavefunctions may influence the precise ordering wavevector, these details do not affect the phenomenological model or the conclusions of the later sections. The corresponding electronic free energy takes the form
\begin{align}\label{eq:electronFreeEnergy}
\mathcal{F}_{\rm ele} =
a(T,n)\abs{\psi_{\vb*{q}}}^2 + b\abs{\psi_{\vb*{q}}}^4,
\end{align}
with the second-order coefficient [see Appendix \ref{app:derivationFreeEnergy} for the derivation]
\begin{align}\label{eq:electronCoefficient}
a(T,n) = \frac{1}{V_{\vb*{q}}} - \chi_{\vb*{q}}(T,n),
\end{align}
where $V_{\vb*{q}}$ denotes the effective electronic interaction at momentum $\vb*{q}$ after integrating out the high-energy degrees of freedom. In this phenomenological description, we neglect the orbital dependence of the interaction $V_{\vb*{q}}$ and treat it as a constant value across all orbitals [see Appendix \ref{app:derivationFreeEnergy} for the detailed model parameters]. As in a standard second-order phase transition, the sign of the coefficient $a(T,n)$ determines whether the electronic order $\psi_{\vb*{q}}$ breaks symmetry, controlled by the relative strength between the interaction and $\chi_{\vb*{q}}(T,n)$. When $\chi_{\vb*{q}}(T,n)$ is enhanced, as it typically is near a VHS or under a favorable nesting condition, and exceeds $1/V_{\vb*{q}}$, the coefficient $a(T,n)$ becomes negative. The system then stabilizes an electronic order that breaks translational and rotational symmetry. Within this framework, we identify such a dominant or subleading density-wave instability as the microscopic origin of the anisotropic stripe-like state recently observed in transport experiments\,\cite{qin2025extreme, zhao2026hysteretical}.


\begin{figure} [!t]
\centering 
\includegraphics[width=8.5cm]{Figure4_lattice_instability_v3.eps}\vspace{-2mm}
\caption{(a) Schematic of the bottom-layer shear mode in rhombohedral graphene. The A and B sublattices are denoted by open and filled circles, respectively, and the bottom layer is displaced with an amplitude modulated in space by wavevector $\vb*{q}$. (b) DFT-calculated energy landscape associated with the shear-mode lattice distortion in rhombohedral hexalayer graphene, showing multiple metastable minima at finite displacement. The solid red curve denotes the corresponding sixth-order phenomenological fit used in our simulations.}
\label{fig: Figure4_lattice_instability} 
\end{figure}

Beyond the spontaneous symmetry breaking discussed above, recent experiments have uncovered pronounced first-order hysteresis in the stripe states. This hysteresis is highly reproducible under repeated scans along both the temperature and doping axes\,\cite{qin2025extreme, zhao2026hysteretical}, strongly suggesting that the system hosts multiple metastable states. An important feature is that the hysteresis extends into doping regimes beyond the strongly correlated flat-band region\,\cite{zhao2026hysteretical}. This observation implies that the charge-order instability itself, which is strongly doping dependent, is unlikely to be solely responsible for the full hysteretic behavior. We therefore propose that the electronic charge order is coupled to an additional lattice degree of freedom, a possibility that is broadly consistent with phenomena observed in graphene-based systems\,\cite{chen2024strong}. In particular, we emphasize that the lattice is intrinsically soft and that its long-wavelength acoustic phonons are susceptible to strong anharmonic effects, owing both to the unstable crystal structure of rhombohedral graphene itself and to interfacial effects\,\cite{singh2025stacking}. 

As a proof of principle, we present the first-principles simulations of the crystal energy as a function of a representative shear mode in rhombohedral hexalayer graphene, where the top five layers are displaced relative to the bottom layer, which is the direction towards the Bernal configuration, as illustrated in Fig.~\ref{fig: Figure4_lattice_instability}(a). In these calculations, we consider a long-wavelength distortion with wavevector $q_x\sim 1/80$. Although this wavevector does not yet match the Fermi-surface nesting wavevector discussed in Fig.~\ref{fig: Figure2_FS_nesting}, because of the limitation imposed by accessible supercell sizes [see Appendix \ref{app:DFT} for details], it nevertheless captures the generic lattice potential that is qualitatively relevant for our phenomenological model. As shown in Fig.~\ref{fig: Figure4_lattice_instability},the simulated energy landscape is strongly anharmonic: the second metastable state has an energy comparable to that of the relaxed equilibrium state, yielding a potential barrier comparable to the harmonic trapping scale. (We also note that the distortion amplitude around $X \sim 0.5\,a$ is still far from the Bernal structure, which would correspond to a substantially lower energy.) As the distortion increases further, the landscape develops multiple local minima. In principle, each of these metastable configurations could support an independent hysteretic response, provided that external conditions are sufficient to overcome the corresponding barriers. 

Because R6G supports many possible lattice-vibration and distortion channels, we do not restrict the phenomenological theory to this specific shear mode. Here, we only consider the lowest-order anharmonic expansion of the lattice potential, which leads to a sixth-order lattice free energy:
\begin{align}\label{eq:latticePotential}
\mathcal{F}_{\rm lat} =
\alpha \abs{X_{\vb*{q}}}^2 + \beta \abs{X_{\vb*{q}}}^4 + \gamma \abs{X_{\vb*{q}}}^6 ,
\end{align}
where $X_{\vb*{q}}$ denotes a generic lattice distortion at wavevector $\vb*{q}_{\rm nesting}$, and the higher-order terms encode the intrinsic anharmonicity. The coefficients $\alpha$, $\beta$, and $\gamma$ are treated as phenomenological parameters controlled by temperature, and their values are chosen to be consistent with both the experimental hysteresis behavior and the DFT energy landscape [see Appendix \ref{app:latticePotential} for details]. In the present work, the potential in Eq.~\eqref{eq:latticePotential} is designed to capture only the nearest metastable state at finite $X_{\vb*{q}}$ [see the dashed line in Fig.~\ref{fig: Figure4_lattice_instability}(b)], which is sufficient for our purposes here. A more microscopic treatment, including multiple metastable structures, is left for future work aimed at predicting electronic control of structure.

\begin{figure}[!t]
\centering
\includegraphics[width=8.5cm]{Figure5_two_order_scape_v3.eps}\vspace{-3mm}
\caption{ (a,b) Free-energy landscape $\mathcal{F}(\psi, X)$ in the space of electronic order $\psi$ and lattice order $X$ without electron-lattice coupling ($g=0$), shown for the cases of (a) no electronic order and (b) finite electronic order. The top and left insets show the corresponding $\mathcal{F}_{\rm ele}$ and $\mathcal{F}_{\rm lat}$, obtained from the horizontal and vertical cuts marked by dashed lines. The cyan circles indicate local minima of the free-energy landscape and thus metastable electron-lattice states. (c,d) Same as in (a,b), but for finite electron-lattice coupling $g$.
}
\label{fig: Figure5_two_order_scape}
\end{figure}


In the absence of electron-lattice coupling, the two-order phenomenological model $\mathcal{F}=\mathcal{F}_{\rm ele}+\mathcal{F}_{\rm lat}$ gives rise to the potential landscape illustrated shown in Figs.~\ref{fig: Figure5_two_order_scape}(a) and (b). When the electronic sector is not ordered, corresponding to $a(T,n)>0$ in Eq.~\eqref{eq:electronFreeEnergy}, the landscape along the electronic-order direction, \textit{i.e.}~the horizontal cut in Fig.~\ref{fig: Figure5_two_order_scape}(a), has a single minimum at $\psi_{\vb*{q}}=0$. At the same time, the lattice sector, \textit{i.e.}~the vertical cut in Fig.~\ref{fig: Figure5_two_order_scape}(a), retains three local minima because of the anharmonic lattice potential. As a result, the full two-order landscape contains three metastable states. Once the electronic order is established, $\mathcal{F}_{\rm ele}$ develops the familiar Mexican-hat structure with two symmetry-breaking minima. Combined with the anharmonic lattice potential, this expands the number of metastable states to six, which are adiabatically connected to the three states in Fig.~\ref{fig: Figure5_two_order_scape}(a). The evolution of the system among these local minima under changing external conditions will determine the anisotropy and hysteresis discussed below.


The leading-order coupling between the stripe CDW and the lattice distortion is linear and momentum conserving, in the form of $\mathcal{F}_{\rm ele\text{-}lat} = g_{\vb*{q}} \psi_{-\vb*{q}} X_{\vb*{q}}$. Once this coupling term $\mathcal{F}_{\rm ele\text{-}lat}$ is included, the energy landscape is modified as illustrated in Fig.~\ref{fig: Figure5_two_order_scape}(c).  When the electronic sector is not spontaneously ordered, the coupling simply locks the symmetry of $\psi_{-\vb*{q}}$ to that of $X_{\vb*{q}}$, producing only a trivial symmetry breaking. As we will show in Sec.~\ref{sec:dopingHysteresis}, this induced symmetry breaking is extremely small and should not be experimentally observable. The situation changes qualitatively once the electronic order becomes sufficiently strong. In that regime, the six distinct metastable states shown in Fig.~\ref{fig: Figure5_two_order_scape}(b) collapse into four, as illustrated in Fig.~\ref{fig: Figure5_two_order_scape}(d). This restructuring of the metastable landscape, including the adiabatic connection between unordered and ordered states, underlies the hysteretic behavior discussed below.

We note that, for a given lattice order $X_\mathbf{q}$, the coupling constant $g_{\vb*{q}}$ is microscopically not a universal constant, but rather  an effective quantity derived from the underlying electron-lattice interaction vertex and, more generally, should be represented as a matrix $g_{\vb*{k}\vb*{q}}^{(mn)}$ that depends on both momentum and band indices. To gain qualitative insight into this structure, we use the band-resolved deformation profile $|\partial E_n/\partial X_{\vb*{q}}|$ as a relative measure. Fig.~\ref{fig: Figure3_shear_mode}(a) shows the distribution of this deformation across different bands, for the bottom-layer shear mode $X_{\vb*{q}}$ illustrated in Fig.~\ref{fig: Figure4_lattice_instability}(a). Within the flat-band regime, the lattice distortion acts predominantly on the first valence band on the hole-doped side. Although the present phenomenological treatment does not attempt to resolve the charge order into multiple sectors, this projection suggests that the effective coupling weakens as the doping moves away from the VHS [also see Fig.~\ref{fig: Figure3_shear_mode}(b)]. 

\begin{figure} [!t]
\centering 
\includegraphics[width=8.5cm]{Figure3_shear_mode.eps} \vspace{-3mm}
\caption{(a) Band-resolved deformation potential $|\partial E_n/\partial X_{\vb*{q}}|$ for the bottom-layer shear-mode lattice distortion shown in Fig.~\ref{fig: Figure4_lattice_instability}(a), plotted along the $k_x$ cut through valley $K$. (b) Deformation-potential profile projected onto the top valence band in the $k_x$-$k_y$ momentum space.}
\label{fig: Figure3_shear_mode} 
\end{figure}




\begin{figure*} [!t]
\centering 
\includegraphics[width=18cm]{Figure6_hyst_v8.eps}\vspace{-3mm}
\caption{(a) Doping dependence of the electronic susceptibility $\chi$ at $T=6.4$\,K. The VHS and $n_{\rm bump}$ are indicated by the shaded regime and dashed line, respectively. The inset shows the corresponding low-energy band structure in the relevant doping range, with the VHS and $n_{\rm bump}$ highlighted. (b) Doping dependence of the corresponding quadratic coefficient $a(T,n)$ under the same conditions. The zoomed-in panel highlights the sign change near the VHS induced by the enhanced susceptibility. (c,d) Forward and backward doping sweeps of (c) the electronic order $\psi$ and (d) the lattice order $X$, showing the hysteresis triggered by the VHS. The non-monotonic bump near $n_{\rm bump}$ during the backward sweep is marked. (e) Hysteresis amplitude $|\psi_{\rm bwd}-\psi_{\rm fwd}|$ in the doping-temperature plane, showing a dome-like structure near $n_{\rm bump}$. The dashed line marks the boundary of the hysteretic region. 
}
\label{fig: Figure6_hyst}
\end{figure*}

\section{Doping-Induced Hysteresis and Anisotropy}\label{sec:dopingHysteresis}

The multi-metastable-state potential landscape of this two-order model governs the evolution of the anisotropic stripe order as external conditions are varied. Because the electronic free energy in Eq.~\eqref{eq:electronFreeEnergy} undergoes spontaneous symmetry breaking when the second-order coefficient $a(n,T)$ becomes negative, the resulting order is highly sensitive to both carrier doping $n$ and temperature $T$. These two parameters can be tuned precisely in experiments.

As discussed above, the DOS is maximized near the VHS associated with the top band, and this regime coincides with the development of Fermi-surface nesting between the two triangular pockets. Together, these two effects strongly enhance the susceptibility $\chi$, as shown in Fig.~\ref{fig: Figure6_hyst}(a). In the presence of a phenomenological Coulomb interaction $V_q$, this singular enhancement drives $a(T,n)$ negative and triggers a large electronic order $\psi_q$ at the nesting wavevector $q$. Such an order parameter should be directly measureable by x-ray scattering or STM microscopy. In transport experiments, this electronic order $\psi_q$ manifests indirectly as by the insulating behavior along the symmetry-breaking direction, while the resistance $R_{xx}$ along the orthogonal direction does not have to increase comparably. This naturally produces the anisotropic behavior recently observed near this VHS\,\cite{zhao2026hysteretical}. Because the electronic order is coupled to the lattice degrees of freedom, the onset of stripe order also induces a finite lattice order [see Fig.~\ref{fig: Figure5_two_order_scape}(d)]. The resulting lattice distortion then relaxes into one of the two low-symmetry metastable states.

When the doping is tuned away from the VHS, the susceptibility $\chi$ rapidly decreases, as shown in Fig.~\ref{fig: Figure6_hyst}(a). Consequently, the second-order coefficient $a(T,n)$ in the electronic free energy quickly becomes positive as the doping is reduced. For a conventional second-order CDW transition, this would imply that the stripe order disappears immediately at the critical doping $n_c$ defined by $a(T,n_c)\sim 0$. Here, however, lattice anharmonicity qualitatively changes that expectation. Once the lattice order has formed, it cannot smoothly return to the high-symmetry state. Instead, as illustrated in Fig.~\ref{fig: Figure5_two_order_scape}(c), the lattice distortion becomes trapped in a metastable low-symmetry configuration even after the electronic instability has disappeared [i.e.~$a(T,n)> 0$]. This metastability allows the lattice order to persist over a broad doping range in which a CDW would otherwise not form spontaneously. Through electron-lattice coupling, the surviving lattice order continues to sustain a strong electronic order, which manifests experimentally as an anisotropic insulating state away from the VHS\,\cite{zhao2026hysteretical}. The resulting hysteretic behavior is therefore a direct consequence of lattice anharmonicity, and it is triggered only when the doping has crossed the VHS; otherwise, it does not emerge, in agreement experimental observations\,\cite{zhao2026hysteretical}.

Although the mechanism underlying this hysteretic behavior is similar to that responsible for the recently observed nonlinear-Hall anomaly in TaIrTe$_4$\,\cite{tang2026bistable}, the anisotropic hysteresis in RMG carries a distinctive signature of its own electronic structure. In the regime where electronic order cannot arise spontaneously [i.e.~$a(T,n)> 0$], the coefficient $a(T,n)$ controls the curvature of the free-energy landscape along the $\psi$ direction in Fig.~\ref{fig: Figure5_two_order_scape}(c). As a result, a smaller value of $a(T,n)$ allows the same lattice order $X$ to induce a stronger electronic order. In other words, the amplitude of the electronic order parameter $\psi$, which can be inferred from $R_{xx}$ in anisotropic transport, continues to encode the underlying electronic susceptibility through Eq.~\eqref{eq:electronCoefficient} even without spontaneous symmetry breaking. As mentioned above, a RMG-specific feature is the secondary enhancement of DOS and susceptibility arising from the band top of the lower band, denoted as $n_{\rm bump}$ in Fig.~\ref{fig: Figure6_hyst}(a), in addition to the dominant VHS. Although this secondary enhancement is not strong enough to drive a spontaneous CDW order, it is nevertheless imprinted in the increase of $\psi$ within the hysteretic regime. The resulting non-monotonic behavior reflects the proximity to an incipient electronic instability at lower filling, highlights the distinctive multi-flat-band structure of RMG, and is consistent with the experimentally measured hysteresis in $R_{xx}$\,\cite{zhao2026hysteretical}.

Because both the anisotropy and the doping hysteresis originate from the electronic susceptibility of the low-energy states, they are naturally sensitive to temperature through the thermal redistribution of electronic occupation. As the temperature rises, the susceptibility $\chi$ near the VHS is rapidly suppressed. This in turn drives the GL coefficient $a(T,n)$ of the CDW order from negative to positive, eliminating the spontaneous CDW instability and causing the anisotropy to disappear at a critical temperature $T_c$. (For the specific set of model parameters used here, we obtain $T_c = $8.7 K; however, the precise value is not essential, since the main point is the general qualitative behavior.) At the same time, increasing temperature also suppresses the susceptibility at $n_{\rm bump}$, corresponding to the top of the second band. Consequently, both the amplitude of the hysteresis, reflected in $R_{xx}$, and its extent along the doping axis decrease as the system approaches $T_c$. Together, these effects produce a dome-like structure in the $T$-$n$ phase diagram.


\section{Magnetic Suppression of the Charge Order}

In our framework, the stripe-like charge order near the VHS is governed by the charge susceptibility $\chi(\mathbf{q})$, which is set primarily by the Fermi-surface nesting condition. A central consequence is that any external perturbation that disrupts this nesting should suppress the susceptibility and thereby weaken the anisotropy. This expectation can be tested through the dependence on magnetic field or on the magnetization of the electronic states, which provides a direct verification of the mechanism described above.


\begin{figure} [!t]
\centering 
\includegraphics[width=8.5cm]{Figure7_magnetic_dep_v8.eps}\vspace{-3mm}
\caption{
(a) Fermi surfaces near the VHS without magnetic field (left) and under a finite magnetic field (right), corresponding to a Zeeman gap of $\Delta E_B \approx 2$\,meV. (b) Real part of the electronic susceptibility, $\mathrm{Re}(\chi)$, as a function of doping and magnetic field (expressed as the Zeeman gap) at $T=6.4$\,K, showing the progressive suppression of the susceptibility with increasing magnetic field.
}
\label{fig: Figure7_magnetic_dep}
\end{figure}
 


An external perpendicular field $B_{\perp}$ affects the low-energy electronic structure through the Zeeman splitting $\Delta E_B = g \mu_B B_{\perp}$, which defines the relevant magnetic-field scale in the electronic energy. Because intrinsic magnetic interactions in RMG\,\cite{huang2023spin}, which are not included in our phenomenological model, can substantially renormalize the effective $g$ factor away from its bare value, we use the Zeeman splitting energy $\Delta E_B$ itself to parametrize the effect of magnetic field. This splitting shifts the up- and down-spin bands in opposite directions by $\Delta E_B/2$. The resulting offset suppresses the electronic susceptibility in two complementary ways. First, the original VHS in the DOS is split into two smaller peaks, thereby lowering the DOS at the VHS. Second, the relative displacement of the spin-resolved Fermi surfaces distorts the original nesting condition, so that the CDW ordering wavevector can no longer efficiently connect the relevant states. As illustrated in Fig.~\ref{fig: Figure7_magnetic_dep}(a), when the up-spin sector is tuned to the VHS and develops a nested Fermi surface, the down-spin sector is already detuned from the VHS and its Fermi surface is no longer nested, which blocks part of the available scattering channels. The combined effect is a suppression of the charge susceptibility $\chi(\mathbf{q})$, together with a broadening along the doping axis, as the magnetic field increases [see Fig.~\ref{fig: Figure7_magnetic_dep}(b)]. Because nesting is strongest near the VHS, magnetic-field suppression of susceptibility is much more pronounced there than near $n_{\rm bump}$.

This reduced susceptibility is then translated into an increase in the quadratic coefficient $a(n)$ of the electronic free energy $\mathcal{F}_{\rm ele}$. Since the magnetic-field effect on the susceptibility at $n_{\rm bump}$ is relatively weak, $a(n_{\rm bump})$ does not vary appreciably with magnetic field. In contrast, the rapid suppression of the $\chi(\mathbf{q})$ near the VHS drives $a(n_{\rm VHS})$ from negative to positive at a fixed temperature [see Fig.~\ref{fig: Figure8_magnetic_dep2}(a)]. Once this occurs, the CDW order can no longer develop near the VHS, and, accordingly, the hysteresis disappears at the corresponding critical magnetic field. 

Fig.~\ref{fig: Figure8_magnetic_dep2}(b) show representative examples at two different magnetic fields for $T=6.4$\,K. Compared with the zero-field hysteresis in Fig.~\ref{fig: Figure6_hyst}(b), a moderate field corresponding to Zeeman splitting $\Delta E_B \leq 1.5$meV does not significantly alter the hysteresis, whereas a larger splitting of $\Delta E_B = 2$meV, above the critical field, suppresses it completely. We summarize the hysteresis amplitude across different $\Delta E_B$ and doping values in Fig.~\ref{fig: Figure7_magnetic_dep}, which reveals a sharp magnetic-field-induced boundary of the hysteretic regime. This sharp transition is consistent with experiment and further supports the underlying CDW origin of the hysteretic anisotropy\,\cite{zhao2026hysteretical}.

In addition to an external magnetic field, spontaneous spin and valley polarization can likewise suppress the CDW instability by acting as an effective internal Zeeman field. In the fully spin- and valley-polarized regime, such as the quarter-metal phase in RMG, only a single spin--valley component remains over a broad energy range near the Fermi level\,\cite{huang2023spin, koh2024correlated, parramartinez2025band, miao2026various}. As a result, the available DOS is reduced to about $25\%$ of its unpolarized value. Because the charge susceptibility scales with the area of the low-energy phase space available for scattering, $\chi(\mathbf{q})$ in Eq.~\eqref{eq:susceptibility} should be suppressed accordingly by approximately the same factor. Given the fact that the charge order remains tunable over a range of doping and magnetic field, the associated CDW instability is unlikely to be very strong even near the VHS. It then follows that, in the spin-valley-polarized regime, the residual charge susceptibility is insufficient to drive the phase transition. Accordingly, both the electronic order and the associated lattice-induced hysteresis disappear once the system becomes polarized. This behavior also provides a way to distinguish the present charge-order scenario from recently proposed metallic Wigner crystal mechanisms, for which polarization or magnetic field would typically favor localization rather than suppress the ordered state\,\cite{dong2026crystals,feng2026self}.

This suppression of a competing order by internal magnetic polarization is closely related to the Jaccarino--Peter effect widely observed in superconductors., where an external magnetic field compensates an internal field and leads to the re-entrance of superconductivity\,\cite{jaccarino1962ultra,meul1984observation,ran2019extreme,helm2024field,yang2026field,vu2026re}. By analogy, when the subleading charge order is suppressed by an internal magnetic field, a similar re-entrance effect should also occur for the charge order. Specifically, if an external magnetic field is applied to the spin-valley-polarized state along the direction opposite to the internal polarization, it can partially cancel the internal field. The compensation condition is reached near the field where the spin-valley polarization flips, which we denote by $B_c$. Around this critical field, the previously suppressed subleading charge order should reappear, giving rise to a re-entrant hysteretic regime analogous to the Jaccarino--Peter effect. This hysteresis is expected to be confined to a narrow window near $B_c$, because further increasing the magnetic field stabilizes the opposite spin-valley polarization and suppresses the charge order again. Such a re-entrance phenomenon near $B_c$ is consistent with the recent experiment\,\cite{zhao2026hysteretical}.


\begin{figure} [!t]
\centering 
\includegraphics[width=8.2cm]{Figure8_magnetic_dep2_v6.eps}
\caption{
(a) Quadratic coefficient $a(T,n)$ of the electronic free energy in the doping-Zeeman-gap plane. (b) Electronic hysteresis amplitude $|\psi_{\mathrm{fwd}}-\psi_{\mathrm{bwd}}|$ in the same parameter space. The hysteretic region collapses at the critical field $\Delta E_B = 1.7$, where $a(T,n)$ remains positive at the VHS. The VHS and $n_{\rm bump}$ are highlighted. (c) Doping hysteresis loops of the electronic order for $\Delta E_B = 0.5$\,meV to $2$\,meV.
}
\label{fig: Figure8_magnetic_dep2}
\end{figure}

\section{Discussions}

The charge-order and phonon-anharmonicity framework proposed in this paper provides a unified explanation for the recently observed anisotropic resistivity \, \cite{qin2025extreme,morissette2025superconductivity}, doping-induced hysteresis, and the incompatibility of these phenomena with magnetic field or polarization \, \cite{zhao2026hysteretical}. Although the present model is phenomenological and does not exclude alternative mechanisms, it is microscopically grounded in two key ingredients: Fermi-surface nesting and a metastable lattice structure. Together, these ingredients suggest a broader physical setting in which similar behavior may arise. More generally, 2D flat-band materials can provide such a setting because strong electronic correlations may coexist with interfacial lattice instability. In this case, if the electron-lattice coupling is sufficiently strong to overcome the barriers between different metastable lattice states, analogous anisotropy and hysteresis should be expected. A representative example is TaIrTe$_4$, where a similar doping-induced stripe order and hysteresis have been observed, with the triggering doping occurring at a VHS and with direct evidence for lattice bistability\,\cite{tang2026bistable}. In graphene-based systems, doping-induced hysteresis has also been reported in bilayer graphene-hBN heterostructures and tetrahedral graphene systems\,\cite{zheng2023electronic, singh2025stacking}. In those systems, the hysteresis has been attributed to lattice metastability either between graphene layers or between graphene and the substrates, although direct evidence remains limited. While these materials may not exhibit the same anisotropic stripe order discussed in this paper, other electronic instabilities, including excitonic order and ferroelectricity, have been suggested to play the dominant role and to couple to the lattice.

More importantly, the experimental observation of the mechanism described above provides an indirect evidence that electron-lattice coupling plays an important role in RMG. This view is consistent with recent evidence for strong polaronic dressing in twisted bilayer graphene, another flat-band 2D platform\,\cite{chen2024strong}. In RMG, multiple competing orders have now been observed, including fractional Chern insulating states, orbital ferromagnetism, anisotropic stripe order, and chiral superconductivity. Their coexistence and competition resemble the intertwined phenomena widely discussed in cuprates\,\cite{fradkin2015colloquium, davis2013concepts, ghiringhelli2012long, arpaia2019dynamical, comin2015symmetry, lee2020spectroscopic}, and likewise raise the question of the origin of superconductivity\,\cite{jiang2019superconductivity, chung2020plaquette}. On the one hand, ferromagnetic fluctuations are believed to favor spin-triplet pairing, although whether a true long-range ordered state can be established remains to be verified by large-scale numerical simulations\,\cite{xu2024coexistence, jiang2021ground}. On the other hand, electron-phonon coupling has already been shown to be sufficient to induce spin-triplet superconductivity in RMG within weak-coupling BCS theory\,\cite{zan2024electron}. These two pairing mechanisms are not mutually exclusive: in fact, large-scale DMRG simulations of strongly correlated models have shown that the cooperation between strong electronic interactions and phonon-mediated pairing can stabilize $p$-wave superconductivity in quarter-filled 1D systems\,\cite{qu2022spin}. This result may be connected to the recent observation that spin-triplet pairing is enhanced when a quasi-1D stripe pattern forms in the normal state\,\cite{qin2025extreme}. However, detailed numerical simulations specific to the RMG setting will still be needed before any firm conclusion can be drawn.

While our theoretical framework is consistent with existing experimental observations, direct experimental confirmation of charge modulation and electron-lattice coupling is still needed. One important direction would be to look for stripe patterns in STM measurements or band folding in nanoARPES measurements, although both types of measurements are generally challenging in devices that require top and bottom gating. Another route is to probe polaronic dressing and phonon anharmonicity through Raman spectra, where they may appear as phonon softening and abrupt changes of phonon branches. A more indirect signature is the temperature dependence of resistivity. In general, linear temperature dependence is often regarded as evidence of phonon effects when the temperature is sufficiently higher than the relevant phonon energy. In RMG, however, this signature may be difficult to identify directly if the phonon anharmonicity is strong, because superlinear behavior arising from anharmonic energy levels can mask the distinction between linear and higher-order temperature dependences. Finally, from the perspective of sample engineering, lattice anharmonicity should be sensitive to the substrate environment. As a result, both the substrate material and the alignment angle are expected to have a strong influence on the observed hysteresis and anisotropy. These considerations suggest several concrete experimental directions that can test, verify, or falsify the present framework.



\section*{Acknowledgements}
We thank Shuhan Ding for insightful discussions. This work is supported by the Air Force Office of Scientific Research Young Investigator Program under grant FA9550-23-1-0153. The simulations presented in this paper used resources of the National Energy Research Scientific Computing Center, a U.S. Department of Energy Office of Science User Facility located at Lawrence Berkeley National Laboratory, operated under Contract No. DE-AC02-05CH11231.









\appendix



\section{Derivation of the Electronic Free-Energy Model}\label{app:derivationFreeEnergy}

In this section, we derive the effective electronic free-energy model
used in the main text from a microscopic interacting electron model.
The low-energy electrons are described by
\begin{equation}
H
=
H_0
+
H_{\mathrm{int}},
\end{equation}
where
\begin{equation}
H_0
=
\sum_{\mathbf k,n,s}
\epsilon_{n,\mathbf k,s}
d^\dagger_{n,\mathbf k,s}
d_{n,\mathbf k,s}
\end{equation}
is the noninteracting Hamiltonian, where $d_{n,\mathbf k,s}$ is the
annihilation operator on a band basis and $s=\uparrow,\downarrow$ labels
spin.

To capture the low-energy density-wave-like ordering tendency at finite
momentum, we consider an effective momentum-conserving interaction in
the density channel,
\begin{equation}
\label{eq:siInteraction}
H_{\mathrm{int}}
=
\frac12
\sum_{\mathbf q}
V_{\mathbf q}
\rho_{-\mathbf q}
\rho_{\mathbf q},
\end{equation}
where $V_{\mathbf q}$ is the effective interaction strength.
The density operator projected to the band basis is
\begin{equation}
\rho_{\mathbf q}
=
\sum_{\mathbf k,m,n,s}
g_{mn}(\mathbf k,\mathbf q)
d^\dagger_{m,\mathbf k+\mathbf q,s}
d_{n,\mathbf k,s},
\end{equation}
with the band form factor
\begin{equation}
g_{mn}(\mathbf k,\mathbf q)
=
\left\langle
u_{m,\mathbf k+\mathbf q}
\middle|
u_{n,\mathbf k}
\right\rangle .
\end{equation}
The form factor encodes the wave-function overlap between the two
electronic states connected by momentum transfer $\mathbf q$.

To derive the effective free energy, we introduce a Hubbard--Stratonovich
field $\psi_{\mathbf q}$ in the same density channel under static mean-field,
\begin{equation}
e^{-\beta H_{\mathrm{int}}}
=
\prod_{\mathbf q}
\int
\mathcal D[\psi_{\mathbf q}]
e^{-S[d,d^\dagger,\psi_{\mathbf q}]},
\end{equation}
where the auxiliary field is 
\begin{equation}
\label{eq:siHSField}
S[d,d^\dagger,\psi_{\mathbf q}]
=
\frac{
|\psi_{\mathbf q}|^2
}{
2V_{\mathbf q}
}
+
\psi_{-\mathbf q}\rho_{\mathbf q},
\end{equation}
where $\psi_{\mathbf q}$ represents the finite-momentum electronic order
parameter.  This form can be understood by integrating out the HS field,
which recovers the original density-density interaction in
Eq.~\eqref{eq:siInteraction}.  Equivalently, at the saddle point one
obtains $\psi_{\mathbf q}\propto -V_{\mathbf q}\rho_{\mathbf q}$,
showing that a nonzero $\psi_{\mathbf q}$ is the same density-wave order
described by $\rho_{\mathbf q}$.  The overall sign is conventional and
can be absorbed into the phase of $\psi_{\mathbf q}$.

After the HS transformation, the electron action is quadratic in the
fermion fields.  The field $\psi_{\mathbf q}$ acts as a static scattering
potential that mixes electronic states with momenta $\mathbf k$ and
$\mathbf k+\mathbf q$.  Integrating out the electronic degrees of freedom
therefore yields
\begin{equation}
e^{-\beta \mathcal{F}_{\mathrm{ele}}[\psi]}
=
\int
\mathcal D[d,d^\dagger]
e^{-S[d,d^\dagger,\psi]} ,
\end{equation}
or, equivalently,
\begin{equation}
\mathcal F_{\mathrm{ele}}[\psi]
=
\sum_{\mathbf q}
\frac{|\psi_{\mathbf q}|^2}{2V_{\mathbf q}}
-
T\,\mathrm{Tr}
\ln
\left(
G_0^{-1}
+
\Sigma_\psi
\right).
\end{equation}
Here $G_0$ is the Green's function of $H_0$, $\Sigma_\psi$ is the
scattering vertex generated by $\psi_{\mathbf q}$, and $\mathrm{Tr}$
denotes a trace over momentum, Matsubara frequency, band, and spin indices.
Expanding the logarithm gives
\begin{equation}\label{siLogExpansion}
-
T\mathrm{Tr}
\ln
\left(
1
+
G_0\Sigma_\psi
\right)
=
-
T
\sum_{\ell=1}^{\infty}
\frac{(-1)^{\ell+1}}{\ell}
\mathrm{Tr}
\left(
G_0\Sigma_\psi
\right)^\ell .
\end{equation}
The linear term vanishes for a finite ordering wavevector because one
insertion of $\Sigma_\psi$ changes the electron momentum by
$\pm\mathbf q$ and cannot close the fermion loop.  The first nonzero term
from the trace-log expansion is therefore quadratic and represents the
bubble.  Keeping only this term, we obtain
\begin{equation}
\label{eq:siQuadraticTrace}
\mathcal F_{\mathrm{ele}}^{(2)}
=
\sum_{\mathbf q}
\frac{|\psi_{\mathbf q}|^2}{2V_{\mathbf q}}
+
\frac{T}{2}
\mathrm{Tr}
\left(
G_0\Sigma_\psi
G_0\Sigma_\psi
\right).
\end{equation}
For a static density-wave order, $\Sigma_\psi$ contains the scattering
vertex
\begin{equation}
\left(\Sigma_\psi\right)_{m,\mathbf k+\mathbf q;n,\mathbf k}
=
\psi_{-\mathbf q}
g_{mn}(\mathbf k,\mathbf q),
\end{equation}
together with its Hermitian-conjugate process.  In spin space the vertex
is diagonal,
\begin{equation}
\left(\Sigma_\psi\right)_{m,\mathbf k+\mathbf q,s;n,\mathbf k,s'}
=
\psi_{-\mathbf q}
g_{mn}(\mathbf k,\mathbf q)
\delta_{ss'} .
\end{equation}

To identify the trace term, we first recall the density response function
in the normal state,
\begin{equation}
\chi(\mathbf q,\tau)
=
-
\left\langle
T_\tau
\rho_{\mathbf q}(\tau)
\rho_{-\mathbf q}(0)
\right\rangle_{0},
\end{equation}
where the expectation value is evaluated using $H_0$.  Substituting the
projected density operator into this expression and applying Wick's
theorem gives the Matsubara-frequency bubble
\begin{align}
\chi(\mathbf q,i\Omega_l)
=&
-
T
\sum_{\omega_n,\mathbf k,m,n,s}
\left|
g_{mn}(\mathbf k,\mathbf q)
\right|^2
\nonumber\\
&\times
G_{0,m,s}(i\omega_n+i\Omega_l,\mathbf k+\mathbf q)
G_{0,n,s}(i\omega_n,\mathbf k).
\end{align}
The spin index is the same on the two Green's functions, reflecting that
only spin-conserving particle-hole excitations contribute to the charge
susceptibility.  Equivalently, $\chi=\chi_\uparrow+\chi_\downarrow$.
Substituting the same vertex into the trace in
Eq.~\eqref{eq:siQuadraticTrace} gives precisely this density bubble.
Therefore, in the static limit relevant for the
mean-field free energy,
\begin{align}\label{eq:siBubbleTrace}
\frac{T}{2}
\mathrm{Tr}
\left(
G_0\Sigma_\psi
G_0\Sigma_\psi
\right)
=&
-
\frac{1}{2}
\sum_{\mathbf q}
\chi(\mathbf q;T,n)
|\psi_{\mathbf q}|^2 ,
\end{align}
where $\chi(\mathbf q;T,n)\equiv \chi(\mathbf q,i\Omega_l=0)$ after the
usual analytic continuation.  The minus sign is the fermion-loop sign in
the bubble definition.  Equivalently, the quadratic part of the free
energy can be written as
\begin{equation}
\mathcal F_{\mathrm{ele}}^{(2)}
=
\sum_{\mathbf q}
\frac{1}{2}
\left(\frac{1}{V_{\mathbf q}}
-
\chi(\mathbf q;T,n)
\right)
|\psi_{\mathbf q}|^2 .
\end{equation}
The factor $1/2$ appears because the full momentum summation counts the
pair $\mathbf q$ and $-\mathbf q$ twice, while
$\psi_{-\mathbf q}=\psi_{\mathbf q}^{*}$ for a real density modulation.
In the main text we keep one independent ordering wavevector, so this
overall factor is absorbed into the convention for the Landau functional.

Including the next nonvanishing term in the trace-log expansion gives the
quartic contribution.  In the main text we keep only its symmetry-allowed
form, $b_{\mathbf q}|\psi_{\mathbf q}|^4$, and take $b_{\mathbf q}>0$ so
that the free energy is stable.  The resulting Ginzburg--Landau free
energy is
\begin{equation}
\mathcal F_{\mathrm{ele}}
=
\sum_{\mathbf q} \left[
a(\mathbf q;T,n)
|\psi_{\mathbf q}|^2
+
b_{\mathbf q}
|\psi_{\mathbf q}|^4
+
\mathcal{O}(|\psi_{\mathbf q}|^6) \right],
\end{equation}
where $a(\mathbf q;T,n) = 1/V_{\vb*{q}} - \chi_{\vb*{q}}(T,n)$.

We now evaluate the static susceptibility entering this coefficient.
The noninteracting Green's function is
\begin{equation}
G_{0,n,s}(i\omega_n,\mathbf k)
=
\frac{1}{i\omega_n-\epsilon_{n,\mathbf k,s}}
\end{equation}
and the Matsubara frequency sum in the bubble is
\begin{align}
&T
\sum_{\omega_n}
G_{0,m,s}(i\omega_n+i\Omega_l,\mathbf k+\mathbf q)
G_{0,n,s}(i\omega_n,\mathbf k)
\nonumber\\
&\hspace{1.0cm}
=
\frac{
f(\epsilon_{n,\mathbf k,s})
-
f(\epsilon_{m,\mathbf k+\mathbf q,s})
}{
i\Omega_l
+
\epsilon_{n,\mathbf k,s}
-
\epsilon_{m,\mathbf k+\mathbf q,s}
}.
\end{align}
Taking the static limit after analytic continuation,
$i\Omega_l\rightarrow \omega+i\eta$ and then $\omega\rightarrow 0$,
therefore leads to the bubble form used in the main text,
\begin{equation}
\chi(\mathbf q;T,n)
=
\sum_{\mathbf k,m,n,s}
\left|
g_{mn}(\mathbf k,\mathbf q)
\right|^2
\frac{
f(\epsilon_{n,\mathbf k,s})
-
f(\epsilon_{m,\mathbf k+\mathbf q,s})
}{
\epsilon_{m,\mathbf k+\mathbf q,s}
-
\epsilon_{n,\mathbf k,s}
+
i\eta
}.
\end{equation}

\section{Model Parameters for the Lattice Potential}\label{app:latticePotential}


\section{Electronic Correction from Electron-Phonon Coupling}
\label{app:ephElectronicCorrection}

We now include the coupling between the electronic density and a static
lattice displacement $X_{\mathbf q}$.  The total Hamiltonian becomes
\begin{equation}
H
=
H_0
+
H_{\mathrm{int}}
+
H_{\mathrm{e\text{-}ph}},
\end{equation}
where $H_{\mathrm{int}}$ is still the density-density interaction in
Eq.~\eqref{eq:siInteraction}.  The electron-phonon coupling is taken to
be in the same density channel,
\begin{equation}
H_{\mathrm{e\text{-}ph}}
=
\sum_{\mathbf q}
g_{\mathbf q}X_{-\mathbf q}\rho_{\mathbf q},
\end{equation}
where $g_{\mathbf q}$ is the electron-phonon coupling constant.  In this
section $X_{\mathbf q}$ is treated as a given lattice coordinate and is
not integrated out.

Since $H_{\mathrm{e\text{-}ph}}$ couples linearly to the same density
operator as the HS field, it can be included in the HS transformation as
an external source.  The interaction and source part is decoupled as
\begin{equation}
e^{-\beta (H_{\mathrm{int}}+H_{\mathrm{e\text{-}ph}})}
=
\prod_{\mathbf q}
\int
\mathcal D[\psi_{\mathbf q}]
e^{-S[d,d^\dagger,\psi_{\mathbf q},X_{\mathbf q}]},
\end{equation}
where the shifted action is
\begin{equation}
S[d,d^\dagger,\psi_{\mathbf q},X_{\mathbf q}]
=
\frac{
\left|
\psi_{\mathbf q}
+
g_{\mathbf q}X_{\mathbf q}
\right|^2
}{
2V_{\mathbf q}
}
+
\psi_{-\mathbf q}\rho_{\mathbf q}.
\end{equation}
This equation is the same HS identity as before, but shifted by the
external source $g_{\mathbf q}X_{\mathbf q}$.  Integrating out the HS
field recovers both the attractive density-density interaction and the
linear electron-phonon coupling.
Use the same expansion in Eq.~\eqref{eq:siLogExpansion} and keep the bubble term, we obtain
\begin{equation}
\label{eq:siQuadraticTrace}
\mathcal F_{\mathrm{ele}}
=
\frac{
\left|
\psi_{\mathbf q}
+
g_{\mathbf q}X_{\mathbf q}
\right|^2
}{
2V_{\mathbf q}
}
+
\frac{T}{2}
\mathrm{Tr}
\left(
G_0\Sigma_\psi
G_0\Sigma_\psi
\right).
\end{equation}

The important point is that the phonon displacement does not enter the
fermionic bubble directly at this order.  After the HS transformation, the
electrons still see the scattering vertex generated by
$\psi_{\mathbf q}$,
\begin{equation}
\left(\Sigma_\psi\right)_{m,\mathbf k+\mathbf q,s;n,\mathbf k,s'}
=
\psi_{-\mathbf q}
g_{mn}(\mathbf k,\mathbf q)
\delta_{ss'} .
\end{equation}
Therefore the trace-log expansion gives the same susceptibility correction
as in Eq.~\eqref{eq:siBubbleTrace}
\begin{align}
\frac{T}{2}
\mathrm{Tr}
\left(
G_0\Sigma_\psi
G_0\Sigma_\psi
\right)
=&
-
\frac{1}{2}
\sum_{\mathbf q}
\chi(\mathbf q;T,n)
|\psi_{\mathbf q}|^2 .
\end{align}
the correction to the Ginzburg--Landau free energy can be written as
\begin{align}
\Delta \mathcal F_{\mathrm{ele}}^{(ph)}[\psi,X]
=
&-
\sum_{\mathbf q}
\frac{1}{2V_{\mathbf q}}
g_{\mathbf q}X_{\mathbf q}\psi_{-\mathbf q}
+
\sum_{\mathbf q}
\frac{|g_{\mathbf q}X_{\mathbf q}|^2}{2V_{\mathbf q}},
\end{align}
where the quadratic coefficient of the lattice order can be absorbed into the lattice free energy model.


\section{First-Principles Simulations}\label{app:DFT}

\bibliography{refs}

\end{document}